\begin{papercard}{gabriele2020}
\cf{Stages} \stagetag{H} (a list of 20 concrete threat scenarios and how plausible and likely users find them), \stagetag{S} (knowledge, attitude and behaviour constructs), \stagetag{R} (survey reporting with non-parametric tests). A CHI 2020 online survey of 212 fitness-tracker users about their security and privacy knowledge, attitudes and behaviours. It concerns wearables and tracker apps rather than mental-health or cycle apps, but its scenario list and sharing-comfort design are directly reusable for the harm stage.

\cf{Question and background} Do users willingly or unknowingly share fitness-tracker data, are they comfortable, and do they understand the implications? \pc{gabriele2020}{1}{Do users willingly or unknowingly share their fitness tracker data?} Motivating incidents: Strava's 2018 heat map revealing military bases, the 2018 MyFitnessPal breach (150 million accounts), employer-sponsored wellness monitoring \pc{gabriele2020}{1}{In 2018, Strava uploaded a heat map of users’ anonymous and aggregate fitness tracking data on their website}. Prior work (reviewed by the authors): users are ill-informed about collection, ownership, sharing, storage and retention; they trust companies; they report the privacy paradox (concern without protective behaviour), though the privacy calculus may explain it better; users rarely read policies; compliance with access requests is poor (Hilts et al.: of 9 companies asked by participants, 5 replied and only Apple fulfilled all requests) \pc{gabriele2020}{3}{only 5 of the 9 companies replied formally, and only one company (Apple) fulfilled all requests}. Identified gap: few studies quantify willingness to share specific data types with specific recipients, and none test scenarios of misuse for fitness trackers. Three results are claimed: lack of knowledge of implications; complex attitudes to risk (fear, likelihood, consequences, confidence in developers); preferences for sharing specific data with specific groups \pc{gabriele2020}{2}{The results demonstrate users’: 1) lack of knowledge of the potential implications associated with the collection of their data}.

\cf{Method} Online survey in January 2019, recruited through Prolific and hosted on Qualtrics. Pre-screen of 508 people; 285 met the criteria (18+, comfortable in English, regular tracker use); 252 completed; 40 excluded for inconsistent demographics between pre-screen and survey or failed attention checks; \textbf{212 valid} \pc{gabriele2020}{4}{Of the 508 participants who completed the pre-screen survey, 285 met the criteria}. The survey (about 20 minutes, 100 questions including 4 attention checks; yes/no/don't know, multiple choice, 5-point Likert, 6 open text) was piloted with 6 participants \pc{gabriele2020}{4}{The final survey has 100 questions, including 4 attention-check questions}; ethics board cleared. Blocks: demographics; confidence in knowledge of data collection and use; \textbf{knowledge of potential threats}: 20 scenarios each answered yes/no/don't know for ``could possibly occur'' and ``would likely occur'' (chosen to force a position; the first scenario: an employer fires an employee after tracker data show late arrival and breaks) \pc{gabriele2020}{4}{participants answered ‘yes’, ‘no’, ‘don’t know’ to whether the scenario “could possibly occur” and whether it “would likely occur”}; attitudes: five concerns (T1 data shared without permission, T2 tampering, T3 account hijacking, T4 data used against you, T5 identity theft after a breach) each rated for fear, likelihood, severity and confidence in developers, based on psychological risk-perception factors (Skirpan et al.); behaviours: protective practices, comfort sharing with 6 recipient groups, and which of 14 Fitbit data types to share with whom. \emph{Analysis}: descriptive statistics and graphics; Friedman rank-sum test with Nemenyi post hoc and Wilcoxon signed-rank for within-subject comparisons \pc{gabriele2020}{4}{Friedman Rank Sum Test, the Nemenyi post-hoc test, and the Wilcoxon Signed-rank Test}; all data self-reported. Sample: 113 female, 98 male, 1 other; ages 18--73 (90\,\% under 45); 55\,\% degree; Fitbit 36\,\%, Apple 17\,\%, Samsung 12\,\%.

\cf{Results: knowledge and scenarios} Participants feel fairly confident about \emph{what} data are collected but less about \emph{how} they are used. For the 20 scenarios the heat map shows many are judged possible but unlikely \pc{gabriele2020}{5}{A heat map of “yes” responses (Figure 2) shows that while participants thought several scenarios were plausible, they also believed many of them unlikely to occur}. Highest possible/likely (\%): mapped running routes revealing movements 94/83; targeted ads for running shoes 91/83; insurance discounts from fitness data 87/71; heart-rate data plus video helps identify a murder suspect 86/60; hacker infers location and activity 85/52; life-insurance cancellation 79/58; police use of tracker data 74/50. Lower: employee fired 71/27; ransomware 70/28; sleep data used against a job applicant 61/21; a user's sexual activity revealed from heart rate and movement 48/19; employer using menstrual and ovulation-cycle data for planning 42/15; adoption agency 41/10; gait identification 32/16; discrimination by race, religion, gender 21/10. The authors note that the four top items had been in the news or resemble common practice.

\cf{Results: attitudes and behaviour} Median fear 3, 2, 2, 3, 3 for T1--T5 (low to mid); likelihood mostly 2; \emph{severity} 4, 4, 4, 4, 5 (``quite severe'') \pc{gabriele2020}{6}{Severity of the consequences – The median responses to Concerns T1 – T5, respectively, were 4, 4, 4, 4, 5}; confidence in developers about the mid-point (3). Interest in more information rose from median 3 to 4 between the beginning and end of the survey (Wilcoxon $W=860.5$, $r=.54$), with the caution that the survey itself primes concern \pc{gabriele2020}{6}{A Wilcoxon Signed-rank Test showed a significant difference in interest}. Behaviour: most do not change default settings, read terms and conditions or read the policy; only about half set sharing preferences \pc{gabriele2020}{6}{They do not change the default security and privacy settings, do not read the terms and conditions, and do not read the privacy policy}. Comfort sharing: highest with the tracker itself and friends, then colleagues, employers, the public, insurance companies, and \emph{lowest with advertisers} (Friedman test significant; Nemenyi pairs mostly different) \pc{gabriele2020}{7}{A Friedman rank sum test revealed a significant difference among the recipients with whom participants were comfortable sharing data}. Data types: gender, height, birthday shareable with many groups; pictures, location and friends list least; sleep and weight graphs less shareable than steps and distance; participants did not distinguish among the three statistical items.

\cf{Interpretation and recommendations} Seven insights, e.g.\ users believe they know what is collected but not how it is used, understanding of threats is limited, concerns are inconsistent (low fear, low likelihood, high severity) \pc{gabriele2020}{7}{Users believe they know what’s collected but are unsure how it’s used}; comfort sharing with the tracker implies unawareness of harm from collection itself \pc{gabriele2020}{9}{Interestingly, they are also comfortable sharing with the fitness tracker itself and by extension the device manufacturer}. Recommendations: awareness education centred on informed consent; granular control over collection and recipients with custom groups; nudges to revisit settings; ``just in time'' consent without punishing users for refusing \pc{gabriele2020}{10}{Efforts towards “just in time” information related to sharing and informed consent should be considered}.

\cf{Limits (authors and card)} Sample not representative (90\,\% aged 18--44, Prolific users in several countries); self-reports may not match behaviour; interviews would add nuance; no ``family'' or physician recipient group \pc{gabriele2020}{10}{Our survey collects self-reported subjective responses and may not reflect participants’ actual behaviours}. Card additions: scenario wording is the authors' own and mixes events that have occurred with hypotheticals, so the possible/likely split partly reflects media exposure; ``likely'' is asked as a yes/no judgement, not a probability; a single wave; ordinal medians without effect sizes for most contrasts; the survey asks about Fitbit-like items, not the in-app consent or policy; no coder or reliability aspects (closed questions; open-text answers are not analysed in the extract).

\cf{Use for this thesis} (1) A model for the user-facing side of stage H: concrete misuse scenarios (insurance, employer, advertiser, law enforcement, inference of sexual activity or menstrual cycle) with ``plausible'' versus ``likely'' ratings; these can seed harm-tree branches and a short list of scenarios to check against what the apps' policies actually permit. (2) The recipient set (tracker/company, friends, colleagues, employers, public, insurers, advertisers) is a vocabulary for transmission-principle recipients in policy coding. (3) The finding that users distrust advertisers more than the public, and are comfortable sharing with the vendor, is evidence that vendor-side collection is under-perceived as a risk. (4) The knowledge/attitude/behaviour split and the self-report limits are a caution when interpreting any thesis claim about what users ``understand''; the thesis analyses materials, not users.
\end{papercard}
